\section*{Bab \ref{ch:b1:nucleophilic-substitution} --- Substitusi Nukleofilik}

\begin{solution}{exo:b1:nucleophilic-substitution:1}
\ce{CH3CH2OH + Br-}; \ce{CH3OCH2CH3 + I-}; \ce{CH3CH2CH2CN + Cl-};
\ce{CH3NH3+ + Br-}. \omterm{def:b1:nucleophilic-substitution:substitution}{Substrat}: haloalkana-haloalkana itu; \omterm{def:b1:electronic-effects:nucleophile}{nukleofil}:
\ce{OH-}, \ce{CH3CH2O-}, \ce{CN-}, \ce{NH3}; \omterm{def:b1:electronic-effects:nucleophile}{gugus pergi}: ion-ion halida.
\end{solution}

\begin{solution}{exo:b1:nucleophilic-substitution:2}
Orde dua: tiga kali; keduanya dibuat setengahnya, dibagi 4. Orde satu
hanya terhadap RY: tidak berubah; setengahnya.
\end{solution}

\begin{solution}{exo:b1:nucleophilic-substitution:3}
SN2 (primer, \omterm{def:b1:electronic-effects:nucleophile}{nukleofil} yang baik, aprotik); SN1 (tersier, air); SN2
(metil, \omterm{def:b1:electronic-effects:nucleophile}{nukleofil} kuat); \omterm{def:b1:nucleophilic-substitution:sn1}{solvolisis} SN1 (sekunder, \omterm{def:b1:electronic-effects:nucleophile}{nukleofil} netral yang
lemah, \omterm{def:b1:intermolecular-forces-solvents:solvent-classes}{pelarut protik}), mungkin disertai sebagian SN2.
\end{solution}

\begin{solution}{exo:b1:nucleophilic-substitution:4}
Sianida mengambil tempat yang berlawanan dengan brom dan ketiga gugus
lainnya terbalik. Pada produknya gugus CN berperingkat pertama, seperti
Br sebelumnya, dengan peringkat lainnya tidak berubah: susunan yang
terbalik mempunyai deskriptor yang berlawanan,
\cip{S}-2-metilbutananitril.
\end{solution}

\begin{solution}{exo:b1:nucleophilic-substitution:5}
1. \omterm{def:b1:electronic-effects:cleavage}{Heterolisis}: panah dari ikatan C--Br ke Br, menghasilkan
\ce{(CH3)3C+} dan \ce{Br-} (lambat). 2. \omterm{def:b1:lewis-resonance-vsepr:lewis}{Sepasang elektron bebas} air
menuju karbon kationik: \ce{(CH3)3C-OH2+}. 3. \omterm{def:b1:lewis-resonance-vsepr:lewis}{Sepasang elektron bebas}
molekul air lain mengambil sebuah proton gugus \ce{OH2+}, dengan pasangan
O--H kembali ke oksigen. Hanya tahap pertama yang lambat yang menentukan
laju; air, pelarutnya, begitu berlebih sehingga konsentrasinya tidak
berubah.
\end{solution}

\begin{solution}{exo:b1:nucleophilic-substitution:6}
Karbon yang membawa Br bersifat primer, tetapi di sebelahnya gugus
tert-butil mengisi ruang di belakang ikatan C--Br: \omterm{def:b1:elementary-steps:energy-profile}{keadaan transisi} SN2
sangat padat. SN1 akan memerlukan \omterm{def:b1:electronic-effects:intermediates}{karbokation} primer, yang terlalu tidak
stabil.
\end{solution}

\begin{solution}{exo:b1:nucleophilic-substitution:7}
Iodida menggantikan iodida melalui SN2, setiap kali dengan inversi:
\cip{S} menjadi \cip{R} sampai keduanya sama banyak, sebuah \omterm{def:b1:stereochemistry-in-depth:enantiomers}{campuran
rasemat}. Setiap substitusi mengubah satu molekul \cip{S} menjadi satu
molekul \cip{R}: rotasinya turun senilai dua molekul, sehingga rotasi
meluruh dua kali lebih cepat daripada laju substitusi \omterm{def:b1:nucleophilic-substitution:substitution}{substratnya}.
\end{solution}

\begin{solution}{exo:b1:nucleophilic-substitution:8}
SN2: bromometana $>$ 1-bromopropana $>$ 2-bromopropana $>$
2-bromo-2-metilpropana (kepadatan ruang). SN1: urutan sebaliknya
(kestabilan \omterm{def:b1:electronic-effects:intermediates}{karbokation}), dengan bromometana dan 1-bromopropana hampir
tidak bereaksi.
\end{solution}

\begin{solution}{exo:b1:nucleophilic-substitution:9}
$x_{\text{inv}} - x_{\text{ret}} = 0.12$ dan $x_{\text{inv}} + x_{\text{ret}} = 1$:
\qty{56}{\percent} terinversi, \qty{44}{\percent} tertahan. Sebagian besar
\omterm{def:b1:nucleophilic-substitution:sn1}{rasemisasi} SN1, dengan sedikit kelebihan inversi: bromida yang pergi
masih menutupi sisinya ketika air datang.
\end{solution}

\begin{solution}{exo:b1:nucleophilic-substitution:10}
$v = k_1k_2[\mathrm{RY}][\mathrm{Nu}]/(k_{-1}[\mathrm{Y^-}] + k_2[\mathrm{Nu}])$
(\cref{prop:b1:nucleophilic-substitution:sn1-law}). Menambahkan \ce{Y-}
memperbesar penyebutnya: lebih banyak \omterm{def:b1:electronic-effects:intermediates}{karbokation} yang kembali menjadi
\omterm{def:b1:nucleophilic-substitution:substitution}{substrat} alih-alih ditangkap, dan lajunya turun. Dapat diabaikan jika
$k_2[\mathrm{Nu}] \gg k_{-1}[\mathrm{Y^-}]$, seperti ketika \omterm{def:b1:electronic-effects:nucleophile}{nukleofilnya}
adalah pelarut.
\end{solution}

\begin{solution}{exo:b1:nucleophilic-substitution:11}
Konsentrasi sama: $-\mathrm{d}c/\mathrm{d}t = kc^2$, $1/c = 1/C_0 + kt$,
$t_{1/2} = 1/(kC_0) = \qty{5.0e5}{s}$ (sekitar enam hari). Dengan
hidroksida berlebih lima puluh kali, $[\ce{OH-}] \approx 50C_0$ tetap:
$c = C_0
e^{-k't}$ dengan $k' = 50kC_0 = \qty{1.0e-4}{s^{-1}}$, $t_{1/2} = \ln 2/k' =
\qty{6.9e3}{s}$, sekitar dua jam.
\end{solution}

\begin{solution}{exo:b1:nucleophilic-substitution:12}
Bromida harus mengusir \ce{OH-}, \omterm{def:b1:predominant-reaction:strong-acid}{basa kuat} dan \omterm{def:b1:electronic-effects:nucleophile}{gugus pergi} yang buruk:
tidak ada reaksi. Dalam HBr pekat, OH terprotonasi menjadi \ce{-OH2+}, yang
\omterm{def:b1:electronic-effects:nucleophile}{gugus perginya} adalah air. Alkohol tersier: air lepas, menghasilkan
\omterm{def:b1:electronic-effects:intermediates}{karbokation}, yang ditangkap oleh \ce{Br-} (SN1). Alkohol primer:
\ce{Br-} menyerang karbon \ce{R-OH2+} dari belakang sementara air lepas
(SN2).
\end{solution}

\begin{solution}{pb:b1:nucleophilic-substitution:1}
\textbf{1.} Setiap molekul yang bereaksi menghasilkan satu \ce{H+} dan satu
\ce{Cl-}; konduktivitas adalah jumlah sumbangan ion-ion, yang sebanding
dengan jumlah ion.
\textbf{2.} $[\mathrm{RCl}]/[\mathrm{RCl}]_0 = (\kappa_\infty -
\kappa)/\kappa_\infty$.
\textbf{3.} $\ln[\mathrm{RCl}] = \ln[\mathrm{RCl}]_0 - kt$: grafikkan
$\ln(\kappa_\infty - \kappa)$ terhadap $t$.
\textbf{4.} 0.693, 0.603, 0.513, 0.423, 0.333, 0.153, $-0.117$, $-0.387$.
\textbf{5.} Ya, kemiringannya $-0.0180$ per menit di seluruh data: orde
satu.
\textbf{6.} Orde apa pun terhadap air: konsentrasinya tetap
(degenerasi orde).
\textbf{7.} $k = 0.0180/60 = \qty{3.0e-4}{s^{-1}}$.
\textbf{8.} $t_{1/2} = \ln 2/k = \qty{2.3e3}{s}$, \qty{38.5}{min}.
\textbf{9.} $\ln 10/k = \qty{7.7e3}{s}$, \qty{128}{min}.
\textbf{10.} $\ln(\kappa_\infty - \kappa)$: 0.693 pada $t = 0$, 0.109
pada \qty{10}{min}, dan seterusnya: kemiringan $-0.0584$ per menit,
$k = \qty{9.7e-4}{s^{-1}}$.
\textbf{11.} $9.7/3.0 = 3.2$.
\textbf{12.} Pada \qty{30}{min}: $2.000(1 - e^{-0.0584 \times 30}) = 1.653$,
terukur 1.654.
\textbf{13.} SN1: \omterm{def:b1:nucleophilic-substitution:substitution}{substrat} tersier, \omterm{def:b1:intermolecular-forces-solvents:solvent-classes}{pelarut protik}, \omterm{def:b1:electronic-effects:nucleophile}{nukleofil} netral.
\textbf{14.} Seperti pada \cref{exo:b1:nucleophilic-substitution:5}, dengan
Cl.
\textbf{15.} Ionisasi yang lambat hanya melibatkan \omterm{def:b1:nucleophilic-substitution:substitution}{substrat}: $v =
k[\mathrm{RCl}]$.
\textbf{16.} Hidroksida hanya dapat berperan pada tahap penangkapan yang
cepat, sesudah ionisasi yang menentukan laju; lajunya tidak bergantung
padanya.
\textbf{17.} \omterm{def:b1:electronic-effects:intermediates}{Karbokation} itu datar dan \omterm{def:b1:stereochemistry-in-depth:stereocentre}{akiral}; air menyerang kedua
sisinya sama banyak: alkohol rasemat.
\textbf{18.} Sawar pertama yang tinggi (ionisasi) menuju \omterm{def:b1:electronic-effects:intermediates}{karbokation} dalam
lembah yang dangkal, lalu sawar rendah menuju penangkapan: tahap pertama
adalah \omterm{def:b1:elementary-steps:rds}{tahap penentu laju}.
\textbf{19.} Zat itu akan membentuk \omterm{def:b1:electronic-effects:intermediates}{karbokation} primer (terlalu tidak
stabil), dan air terlalu lemah sebagai \omterm{def:b1:electronic-effects:nucleophile}{nukleofil} untuk SN2 yang cepat.
\textbf{20.} $k = A\,e^{-E_a/RT}$.
\textbf{21.} $E_a = R\ln(k_2/k_1)/(1/T_1 - 1/T_2)$.
\textbf{22.} $E_a = 8.314 \times \ln(9.74/3.00)/(1/298.15 - 1/308.15) =
\qty{9.0e4}{J/mol}$.
\textbf{23.} Tinggi sawar ionisasi, yaitu \omterm{def:b1:elementary-steps:rds}{tahap penentu laju}.
\textbf{24.} $k = 3.0 \times 10^{-4} \exp\bigl(\frac{90\,000}{8.314}
(\frac{1}{298.15} - \frac{1}{318.15})\bigr) = \qty{2.9e-3}{s^{-1}}$.
\textbf{25.} \textbf{$E_a = \qty{90}{kJ/mol}$}.
\end{solution}
